\documentclass[UTF8,12pt]{ctexart}
\usepackage[a4paper,margin=24mm]{geometry}
\usepackage{amsmath,amssymb,bm,graphicx,adjustbox}
\newcommand{\fitmath}[1]{\begin{adjustbox}{max width=\linewidth}$\displaystyle #1$\end{adjustbox}}
\setlength{\parindent}{0pt}
\setlength{\parskip}{.7em}
\sloppy
\allowdisplaybreaks
\begin{document}
\section*{2010 年考研数学三 · 参考答案}
\subsection*{原 PDF 第 50 页}

\subsection*{2010年真题参考答案}

\subsection*{一、选择题}

（1）C。 （2）A。 （3）B。 （4）C。 （5）A。 （6）D。 （7）C。 （8）A。


\subsection*{二、填空题}

（9）\(\displaystyle -1\)。 （10）\(\displaystyle \dfrac{\displaystyle \pi^2}{\displaystyle 4}\)。 （11）\(\displaystyle pe^{(p^3-1)/3}\)。 （12）\(\displaystyle 3\)。 （13）\(\displaystyle 3\)。 （14）\(\displaystyle \sigma^2+\mu^2\)。


\subsection*{三、解答题}

（15）\(\displaystyle e^{-1}\)。


（16）\(\displaystyle \dfrac{\displaystyle 14}{\displaystyle 15}\)。


（17）最大值为 \(\displaystyle 5\sqrt5\)，最小值为 \(\displaystyle -5\sqrt5\)。


（18）（I）\(\displaystyle \int_0^1|\ln t|[\ln(1+t)]^n\,dt<\displaystyle \int_0^1t^n|\ln t|\,dt\; (n=1,\allowbreak 2,\allowbreak \ldots)\)。
（II）\(\displaystyle \lim_{n\to\infty}u_n=0\)。


（19）证明略。


（20）（I）\(\displaystyle \lambda=-1,\allowbreak a=-2\)。
（II）\(\displaystyle k(1,\allowbreak 0,\allowbreak 1)^{\mathrm T}+\left(\dfrac{\displaystyle 3}{\displaystyle 2},\allowbreak -\dfrac{\displaystyle 1}{\displaystyle 2},\allowbreak 0\right)^{\mathrm T}\) 为 \(\displaystyle A\boldsymbol x=\boldsymbol b\) 的通解，其中 \(\displaystyle k\) 为任意常数。


（21）\(\displaystyle a=-1\)，\(\displaystyle Q=\)


\[\fitmath{\displaystyle Q=\begin{pmatrix}\dfrac{\displaystyle 1}{\displaystyle \sqrt6}&\dfrac{\displaystyle 1}{\displaystyle \sqrt2}&\dfrac{\displaystyle 1}{\displaystyle \sqrt3}\\\dfrac{\displaystyle 2}{\displaystyle \sqrt6}&0&-\dfrac{\displaystyle 1}{\displaystyle \sqrt3}\\\dfrac{\displaystyle 1}{\displaystyle \sqrt6}&-\dfrac{\displaystyle 1}{\displaystyle \sqrt2}&\dfrac{\displaystyle 1}{\displaystyle \sqrt3}\end{pmatrix},\quad Q^{\mathrm T}AQ=\begin{pmatrix}2&0&0\\0&-4&0\\0&0&5\end{pmatrix}.}\]


（22）\(\displaystyle A=\dfrac{\displaystyle 1}{\displaystyle \pi},\allowbreak \quad f_{Y\mid X}(y\mid x)=\dfrac{\displaystyle 1}{\displaystyle \sqrt\pi}e^{-(x-y)^2},\allowbreak \ y\in(-\infty,\allowbreak +\infty)\)。


（23）（I）随机变量 \(\displaystyle (X,\allowbreak Y)\) 的概率分布为


\[\fitmath{\displaystyle \begin{array}{c|ccc}X\backslash Y&0&1&2\\\hline0&\dfrac{\displaystyle 1}{\displaystyle 5}&\dfrac{\displaystyle 2}{\displaystyle 5}&\dfrac{\displaystyle 1}{\displaystyle 15}\\1&\dfrac{\displaystyle 1}{\displaystyle 5}&\dfrac{\displaystyle 2}{\displaystyle 15}&0\end{array}}\]


（II）\(\displaystyle -\dfrac{\displaystyle 4}{\displaystyle 45}\)。


\clearpage
\end{document}
